\documentclass[fleqn]{beamer}

\mode<presentation> {
\usetheme{Madrid}
\usecolortheme{dolphin}
}

\usepackage[utf8]{inputenc}
\usepackage[brazilian]{babel}
\usepackage{amsmath}
\usepackage{amssymb}
\usepackage{booktabs}
\usepackage{listings}
\usepackage{xcolor}

\lstset{
	basicstyle=\ttfamily\scriptsize,
	breaklines=true,
	keywordstyle=\color{blue!60!black},
	commentstyle=\color{gray},
	stringstyle=\color{green!45!black},
	showstringspaces=false,
	columns=fullflexible,
	frame=single,
	framesep=4pt,
	aboveskip=6pt,
	belowskip=2pt,
}

\title[Aula 4 -- Localização de Facilidades]{Localização de Facilidades}
\subtitle{Estudo de caso e enunciado da tarefa}
\author[Santi]{Prof. \'Everton Santi}
\institute[UFRN/ECT]{ECT2524 -- Introdução à Otimização \\ Escola de Ciências e Tecnologia -- UFRN}
\date{Aula 4}

\begin{document}

\begin{frame}
	\titlepage
\end{frame}

\begin{frame}
	\frametitle{Sumário}
	\tableofcontents
\end{frame}

%------------------------------------------------------------
\section{Recapitulando}
%------------------------------------------------------------

\begin{frame}
	\frametitle{Da Aula 3}
	\begin{itemize}
		\item Chamamos o AMPL de dentro do Python (\texttt{amplpy}):
		\begin{itemize}
			\item \texttt{ler\_dados} lê a instância de um \texttt{.txt};
			\item o script injeta os parâmetros no modelo e resolve;
			\item um laço percorre várias instâncias e monta uma tabela
				de resultados (objetivo, LB, GAP, tempo, status);
		\end{itemize}
		\item Estudo de caso: o Problema de Transbordo (fábricas
			$\rightarrow$ centros $\rightarrow$ clientes);
		\item Hoje: um segundo estudo de caso, que será a \textbf{tarefa}
			-- mesmos passos, problema diferente.
	\end{itemize}
\end{frame}

%------------------------------------------------------------
\section{O problema de localização de facilidades}
%------------------------------------------------------------

\begin{frame}
	\frametitle{O contexto}
	\begin{itemize}
		\item Um órgão aplica um \textbf{exame nacional} e precisa
			organizar, em cada cidade, onde a prova acontece;
		\item Existe uma lista de \textbf{locais candidatos} -- escolas,
			faculdades, centros de eventos;
		\item Cada candidato mora em algum ponto da cidade e será alocado
			a \textbf{um} local para fazer a prova;
		\item Decisão anual: quais locais usar e como distribuir os
			candidatos entre eles.
	\end{itemize}
\end{frame}

\begin{frame}
	\frametitle{Elementos do problema}
	\begin{itemize}
		\item \textbf{Candidatos}: conjunto de pessoas a alocar (sabe-se
			onde cada uma mora);
		\item \textbf{Locais candidatos}: cada um com uma
			\textbf{capacidade} -- número máximo de candidatos que
			comporta;
		\item \textbf{Limite de locais}: o número de equipes de aplicação
			disponíveis fixa quantos locais podem funcionar no dia
			(chamamos esse número de $p$);
		\item \textbf{Distâncias}: para cada par candidato--local, a
			distância que o candidato percorreria se fosse alocado ali.
	\end{itemize}
\end{frame}

\begin{frame}
	\frametitle{O que precisa ser decidido}
	\begin{itemize}
		\item Quais \textbf{locais abrir} (no máximo $p$);
		\item Para cada candidato, \textbf{em qual local aberto} ele faz a
			prova;
		\item Regras que a solução tem que respeitar:
		\begin{itemize}
			\item cada candidato é alocado a exatamente um local;
			\item um candidato só pode ser alocado a um local que esteja
				\textbf{aberto};
			\item nenhum local recebe mais candidatos do que sua
				\textbf{capacidade};
			\item no máximo $p$ locais são abertos;
		\end{itemize}
		\item \textbf{Objetivo}: fazer os candidatos se deslocarem pouco
			-- mas há duas leituras disso.
	\end{itemize}
\end{frame}

\begin{frame}
	\frametitle{Duas visões do objetivo}
	\begin{itemize}
		\item \textbf{Soma das distâncias} -- minimizar o total (ou a
			média) percorrido por todos os candidatos;
		\begin{itemize}
			\item favorece o resultado \textbf{agregado}: bom para a
				maioria, pode deixar poucos candidatos muito longe;
			\item na literatura: problema da \textbf{$p$-mediana};
		\end{itemize}
		\item \textbf{Maior distância} -- minimizar a distância do
			candidato \textbf{mais prejudicado};
		\begin{itemize}
			\item foca no \textbf{pior caso}: ninguém fica longe demais,
				ao custo de uma soma total maior;
			\item na literatura: problema do \textbf{$p$-centro};
		\end{itemize}
		\item Na tarefa vocês implementam \textbf{as duas} e comparam
			instância a instância -- não há escolha a priori.
	\end{itemize}
\end{frame}

\begin{frame}
	\frametitle{Comparar as duas soluções}
	\begin{itemize}
		\item Os objetivos vivem em escalas diferentes (soma
			\emph{vs.}\ máximo) -- comparar os valores ótimos direto não
			diz muito;
		\item O que se compara é o \textbf{deslocamento dos candidatos}
			em cada solução:
		\begin{itemize}
			\item maior distância e menor distância percorrida;
			\item distância média;
			\item desvio padrão (o quão desigual é a solução);
		\end{itemize}
		\item Mais os indicadores de resolução, com \textbf{5 min} de
			limite: LB, UB, GAP, tempo, status.
	\end{itemize}
\end{frame}

\begin{frame}
	\frametitle{Onde mais esse problema aparece}
	\begin{itemize}
		\item Escolher poucos pontos para atender muita gente espalhada:
		\begin{itemize}
			\item postos de saúde ou pronto-atendimento numa região;
			\item centros de distribuição de uma varejista;
			\item antenas de telefonia, estações de recarga;
			\item escolas, bibliotecas, pontos de coleta;
		\end{itemize}
		\item Sempre a mesma estrutura: um número limitado de
			instalações, demanda distribuída no espaço, custo ligado à
			distância.
	\end{itemize}
\end{frame}

%------------------------------------------------------------
\section{A tarefa}
%------------------------------------------------------------

\begin{frame}
	\frametitle{O que entregar}
	\begin{itemize}
		\item Sigam os \textbf{mesmos três passos} da Aula 3:
		\begin{enumerate}
			\item \textbf{Modelar}: parâmetros, variáveis e restrições --
				comuns às duas formulações -- e as \textbf{duas} funções
				objetivo (a do $p$-centro linearizada com variável
				auxiliar);
			\item \textbf{Modelo em AMPL} (\texttt{facilidades.mod}) com
				as duas \texttt{minimize};
			\item \textbf{Código Python}: \texttt{ler\_dados} + script
				\texttt{amplpy} que resolve as \textbf{três instâncias}
				nas \textbf{duas formulações} (6 execuções) e monta a
				tabela comparativa;
		\end{enumerate}
		\item Reaproveitem a colinha do \texttt{amplpy} da Aula 3 (passar
			escalar, vetor, matriz; configurar o solver; ler objetivo,
			LB, GAP, tempo, status; ler os valores das variáveis para
			reconstruir a alocação).
	\end{itemize}
\end{frame}

\begin{frame}[fragile]
	\frametitle{Formato dos arquivos de instância}
	\begin{lstlisting}
10 5 3
2 3 5 1 4
62.70 111.19 59.65 97.79 4.73
9.23 43.95 39.23 108.70 29.87
...
	\end{lstlisting}
	\begin{itemize}
		\item linha 1: \texttt{n\_candidatos  n\_locais  p};
		\item linha 2: capacidade de cada local ($n\_locais$ valores);
		\item próximas $n\_candidatos$ linhas: matriz de distâncias
			candidato $\times$ local;
		\item valores separados por espaço.
	\end{itemize}
\end{frame}

\begin{frame}
	\frametitle{As três instâncias}
	\begin{center}
	\begin{tabular}{cccc}
		\toprule
		Instância & Candidatos & Locais & $p$ \\
		\midrule
		1 & 10 & 5 & 3 \\
		2 & 100 & 10 & 4 \\
		3 & 5000 & 50 & 30 \\
		\bottomrule
	\end{tabular}
	\end{center}
	\begin{itemize}
		\item A instância 3 é bem maior: parte da tarefa é observar se ela
			é resolvida na otimalidade dentro dos 5 min de limite, ou se
			sobra GAP -- e se isso muda entre p-mediana e p-centro;
		\item Enunciado completo, perguntas de reflexão e os três
			\texttt{.txt} estão na página da aula, no site do curso.
	\end{itemize}
\end{frame}

%------------------------------------------------------------
\section{Referências}
%------------------------------------------------------------

\begin{frame}
	\frametitle{Referências}
	\begin{itemize}
		\item AMPL: \url{https://ampl.com};
		\item Documentação do \texttt{amplpy}: \url{https://amplpy.ampl.com};
		\item Daskin, M. S. \emph{Network and Discrete Location}. 2.\,ed.
			Wiley, 2013 (capítulos de $p$-mediana e $p$-centro).
	\end{itemize}
\end{frame}

\end{document}
